\section{Introduction}
\label{sec:introduction}

How does a language model ``know'' that a washing machine is a household appliance?
GPT-2 Small has 768 dimensions in its residual stream, yet English contains millions of referenceable compound concepts---washing machine, coffee table, red herring, guinea pig.
Even the superposition hypothesis~\citep{elhage2022superposition}, which allows models to store ${\sim}10\times$ more features than dimensions through near-orthogonal packing, cannot accommodate this many dedicated compound directions.
Something must give: either compounds share directions with their components, or the model uses a different mechanism entirely.

{\bf We ask: are compound concepts stored as dedicated representations, or are they composed from their parts?}
Three possibilities exist.
First, the model may develop \emph{holistic} representations---unique directions for ``washing machine'' that are distinct from both ``washing'' and ``machine.''
Second, the model may rely on \emph{compositional construction}, where seeing ``washing'' contextually modifies the ``machine'' representation and makes ``machine'' overwhelmingly likely as the next token, with no unified compound direction needed.
Third, a \emph{hybrid} mechanism may operate, where subtle compound-specific features coexist with largely compositional representations.

Prior work has studied superposition for individual features~\citep{elhage2022superposition}, feature absorption in sparse autoencoders~\citep{chanin2024absorption}, and compound noun semantics in BERT-family models~\citep{ormerod2024kitchen,miletic2023systematic}.
However, no study has directly measured how compound concepts are decomposed by SAEs in autoregressive models, or whether next-token prediction serves as the primary mechanism for compound concept ``storage.''
\citet{park2024linear} tested the linear representation hypothesis on 27 single-token concepts but never multi-token compounds.
\citet{merullo2023language} showed that LLMs implement vector arithmetic for relational tasks, but did not examine whether compound nouns follow the same compositional pattern.

We address this gap with three complementary experiments on \gpttwo (124M parameters, 12 layers, 768 dimensions), examining 21 compound concepts that span the full compositionality spectrum.
Our experiments measure (1) residual stream cosine similarity between compound and component representations across all layers, (2) SAE feature overlap between compounds and their parts using pre-trained SAEs with 24,576 features per layer, and (3) next-token prediction priming---how strongly a modifier like ``washing'' shifts the probability distribution toward its compound head ``machine.''

Our results support the hybrid hypothesis, with composition as the dominant mechanism.
The compound direction is nearly identical to the head noun direction in the residual stream (cosine similarity 0.954), yet SAE decomposition reveals that 56.3\% of active features are unique to the compound context.
The primary mechanism is next-token prediction priming: after ``The washing,'' GPT-2 assigns 47.1\% probability to ``machine'' as the top prediction---a $4{,}221\times$ increase over baseline.
Idiomaticity predicts the degree of unique representation: ``red herring'' activates 82.1\% unique SAE features while ``kitchen chair'' activates only 46.9\%.

In summary, we make the following contributions:
\begin{itemize}[leftmargin=*,itemsep=0pt,topsep=0pt]
    \item We conduct the first systematic comparison of compound concept representations across three levels of analysis---residual stream geometry, SAE feature decomposition, and next-token prediction---revealing that these levels tell fundamentally different stories about compound storage.
    \item We show that cosine similarity is a blunt instrument for detecting compound-specific processing: despite 0.954 mean cosine similarity between compounds and head nouns, SAE features reveal 56.3\% unique activation, and this unique fraction correlates strongly with linguistic compositionality (Mann-Whitney $p = 0.0005$, Cohen's $d = 2.54$).
    \item We demonstrate that next-token prediction priming is the dominant mechanism for compound concept construction, with priming ratios spanning five orders of magnitude (from $1.3\times$ for ``office desk'' to $28{,}037\times$ for ``guinea pig''), providing a quantitative taxonomy of how ``frozen'' different compounds are in the model's predictions.
\end{itemize}
